% ---------------------------------------------------------------
\section{Comparison with Base Papers and State of the Art}
\label{sec:comparison}

The research problem was framed in Phase~1 around three base papers.
DriftLens~\cite{greco2025driftlens} motivated the latent-space detection layer, the
domino drift effect (DDE)~\cite{szucs2025dde} motivated feature localization, and
proactive Hoeffding trees~\cite{guerrerocano2026proactive} motivated the adaptation
experiments. Average-confidence and ATC estimators~\cite{garg2022atc,kivimaki2025confidence}
supply the impact baselines.

\subsection{Summary of the Base Papers}
Table~\ref{tab:basepapers} summarises each base paper. Its ``Reported results'' column
is quoted from the original papers and marked with an asterisk.

\begin{table*}[t]
\centering
\caption{Base papers. Numbers marked * are quoted from the original publications and were not reproduced in this work.}
\label{tab:basepapers}
\setlength{\tabcolsep}{3pt}
\footnotesize
\begin{tabular}{L{2.0cm}L{2.3cm}L{2.7cm}L{2.4cm}L{1.9cm}L{2.8cm}L{2.3cm}}
\toprule
Base paper & Problem & Method & Data & Metrics & Reported results* & Key limitation \\
\midrule
Greco et al., 2025~\cite{greco2025driftlens} (DriftLens) & real-time unsupervised drift detection for deep classifiers & per-label Gaussian models of PCA-reduced embeddings; Fréchet drift distance; offline thresholds & 17 use cases on text, image and audio (e.g.\ AG News, STL-10, Common Voice) & detection accuracy, harmonic drift detection, runtime, Spearman correlation & better detection in 15/17 use cases*; at least $5\times$ faster*; Spearman 0.88 / 0.99 / 0.85 for sudden / incremental / recurrent drift* & Gaussian, first two moments; unbounded distance; balanced labels; no link to accuracy loss \\
Sz\H{u}cs \& N\'emeth, 2025~\cite{szucs2025dde} (DDE) & localizing feature drift in high dimension without monitoring every feature & drift-probability curve over correlation thresholds; unmonitored features scored from correlated drifting ones & CICIDS, Covertype, Insects, Heartbeats & recall, precision, F1, AUC, runtime vs monitored fraction & Heartbeats F1 0.8808 (all 280 features monitored) and 0.4373 (8 monitored)*; Covertype F1 1.0000 with 16 features* & needs an initial drift and correlated co-drift; flat lift when shifts are induced independently \\
Guerrero Cano et al., 2026~\cite{guerrerocano2026proactive} (PHT) & anticipating incremental concept drift & four proactive Hoeffding-tree variants that move decision boundaries ahead of drift & synthetic RBF streams; 7 real streams incl.\ Electricity and NOAA Weather & prequential accuracy, degradation, recovery time, memory & PHT-MC 0.9574 vs VFDT 0.9406 prequential accuracy on synthetic streams*; advantage ``significantly diminished'' on real streams* & assumes predictable trajectories; prequential (immediate) labels; incremental drift only \\
\bottomrule
\end{tabular}
\end{table*}

\subsection{Quantitative Comparison}
% generated by scripts/build_tables.mjs from results/phase3/dependence_extension_20261116_final_v2/paired_contrasts.csv, results/phase3/v2_final/paired_estimator_bootstrap.csv, results/phase3/v2_final/impact_summary.csv, results/phase3/noaa_weather_1949_1999_v1/evaluator/batch_metrics.csv, results/phase3/adaptation_final/paired_bootstrap.csv
\begin{table*}[t]
\centering
\caption{Quantitative comparison and improvement analysis against each base paper. \emph{Published*}: results quoted from the base paper on its own data and protocol (not reproduced, and not comparable with the other columns). \emph{Base-style}: the base method's re-implemented, style-level counterpart, measured with \idm{} on identical streams. Brackets: 95\% paired seed-bootstrap intervals; n.s.: not significant. Rel.: change relative to the base-style value; for errors a negative change is an improvement. The DriftLens rows use the fresh seeds of the dependence study (930100--930129), so the KS/PCA rate differs from the 8/30 of Table~\ref{tab:core} (seeds 100--129).}
\label{tab:improvement}
\setlength{\tabcolsep}{3pt}
\footnotesize
\begin{tabular}{L{1.7cm}L{2.7cm}L{2.3cm}L{2.5cm}ccL{2.3cm}c}
\toprule
Base paper & Published* & Base-style component & Metric (same streams) & Base-style & \idm{} & Absolute change & Rel. \\
\midrule
DriftLens~\cite{greco2025driftlens} & best detection in 15/17 deep use cases*; Spearman up to 0.99* & frozen-latent (PCA) + feature KS & C6 onset detection (30 seeds) & 53.3\% & 100.0\% & +46.7 pp [30.0, 63.3] & +87.5\% \\
 & & & stationary alarms / 24 batches & 0.83 & 2.37 & +1.53 (cost) &+184\% \\
ATC~\cite{garg2022atc} / AC~\cite{kivimaki2025confidence} & estimators of accuracy under shift (no comparable number) & ATC-style threshold estimator & accuracy MAE, C6 & 0.074 & 0.019 & -0.0547 [-0.0599, -0.0494] & -73.8\% \\
 & & & accuracy MAE, NOAA &0.067 & 0.044 & -0.0223 & -33.4\% \\
 & & & accuracy MAE, C3 &0.304 & 0.303 & -0.0011 [-0.0033, 0.0010], n.s. & -0.4\% \\
DDE~\cite{szucs2025dde} & Covertype F1 1.0000 (full monitoring)*; Heartbeats F1 0.4373--0.6833 at 3\% monitoring* & co-drift inference of unmonitored features & -- (different monitoring budget) & -- & -- & not compared & -- \\
PHT~\cite{guerrerocano2026proactive} & PHT-MC 0.9574 vs VFDT 0.9406 prequential accuracy* & refit on released labels & C3 post-onset accuracy, delay 0 & 0.279 & 0.278 (frozen) & $-0.02$ pp [$-0.11$, 0.07], n.s. & $-0.1$\% \\
\bottomrule
\end{tabular}
\end{table*}

No base paper addresses the target of this study: verified harm of a frozen tabular
model under delayed labels. A same-dataset, same-metric comparison with their
published numbers is therefore not possible. Table~\ref{tab:improvement} keeps the two
kinds of evidence apart. The \emph{Published*} column quotes each base paper on its own
data. The remaining columns compare a re-implemented, style-level counterpart of each
base method with \idm{} on identical streams, seeds and partitions. Every number
without * was measured in this work.

\subsection{Methodological Comparison}
\idm{} differs from the base papers in its objective, its timing model and its
evaluation target. DriftLens and DDE score the \emph{input} distribution and are
evaluated on detecting or localizing injected drift. \idm{} adds an impact channel that
estimates batch accuracy and a harm-concern score. It scores every channel against
verified harm, and it enforces that the evaluator's truth reaches the monitor only
after the release delay. This makes visible a failure that input-only evaluations
cannot detect: a change in $P(Y\mid X)$ with an unchanged $P(X)$ (Section~\ref{sec:res-conditional}).
PHT adapts a learner with immediate labels. \idm{} keeps the model frozen and treats
adaptation as a separate, budget-matched decision. This separation exposed a confound
in which an apparent advantage of about 12 points for local over global refitting (12.45
versus 0.37 points over the frozen model) disappeared once the update clocks and label windows were matched (Section~\ref{sec:ablation}). Unlike DriftLens's
embedding model, the latent test in \idm{} is a three-component PCA, because the
autoencoder proposed in Phase~2 was not implemented.

\subsection{Improvement Analysis}
The right-hand columns of Table~\ref{tab:improvement} give the absolute and relative
change against each base method's re-implemented counterpart.
\emph{DriftLens-style latent detection.} Adding the correlation statistic raised C6
onset detection from 53.3\% to 100\% of seeds (+46.7 points [30.0, 63.3], +87.5\%).
The gain comes from dependence evidence that marginal and PCA-coordinate tests cannot
see, and it costs 1.53 additional stationary alarms per 24 batches (+184\%).
\emph{ATC/AC estimators.} Calibrated average confidence lowered estimation error by
73.8\% relative to the ATC-style estimator on C6 ($-0.0547$ [$-0.0599$, $-0.0494$]) and
by 33.4\% on NOAA. The gain comes from using every calibrated probability rather than
a single confidence threshold. Under conditional-only change (C3) the two estimators
fail equally, and the difference is not significant. \emph{DDE.} No improvement is
claimed. DDE infers drift in features that are \emph{not} monitored, whereas \idm{}
tests every feature. A style-level co-drift proxy on nearly uncorrelated synthetic
features scored close to random (AP 0.12--0.30), but this reflects the absence of
co-drift in those streams rather than a weakness of DDE.
\emph{PHT-style adaptation.} Refitting on released labels did not change post-onset
accuracy ($\pm$0.1 points, not significant), so the frozen monitor neither gains nor
loses against an adaptive learner here. Two further internal gains complete the
picture. Delayed recalibration on UCI cut MAE by 68.9\% at delay 0 (48.0\% at a
two-cohort delay). Monitor features cut the forecaster's Brier score by 27.9\%
relative to time alone.

\subsection{Fairness and Limitations of the Comparison}
Four limitations bound this comparison. First, DriftLens and DDE were not reproduced:
the latent test and the co-drift proxy are simplified ``style'' components, and
labelling them as the original methods would be incorrect. Second, the base papers
use different data (deep unstructured inputs, high-dimensional correlated features,
synthetic incremental streams), different preprocessing and immediate labels.
Third, the ELEC2 and NOAA streams used here are re-ordered or re-derived from their
sources and are not the benchmark files used by PHT. Fourth, the methods in this work
do \emph{not} outperform the alternatives in several settings: the dependence statistic
nearly triples stationary alarms when combined with KS, delayed recalibration did not help on ELEC2, and the
concern rule has lower recall than the point rule and the ATC-style estimator on C6
(Section~\ref{sec:ablation}).
